
%Objetivos
\chapter{Objetivos} \label{objetivos}
El objetivo general de este proyecto es el diseño y desarrollo de una aplicación web para la gestión eficiente de documentos en una organización, facilitando conocer su estado y si han sido almacenados o no. Además, mediante el uso de filtros avanzados, los usuarios podrán identificar y gestionar de manera eficiente los metadatos de estos documentos, mejorando así la integridad y el control de la información. 
La aplicación también será escalable y segura, garantizando su adecuado funcionamiento en diversos entornos y proporcionando una herramienta útil para la gestión de bases de datos.

Los objetivos específicos del proyecto son:
\begin{enumerate}
    \item Desarrollar una interfaz web altamente usable. El diseño de la interfaz será intuitivo y accesible, permitiendo que cualquier usuario, independientemente de su nivel de conocimientos técnicos, pueda interactuar con la aplicación sin dificultad. Esto facilitará su adopción en empresas de distintos tamaños y sectores.

    \item Implementar un sistema basado en arquitectura cliente-servidor para ordenador. La aplicación será accesible a través de un navegador web, sin necesidad de instalación en dispositivos locales. Funcionará bajo una arquitectura cliente-servidor, donde el servidor aloja y gestiona los datos, mientras que los usuarios interactúan con la aplicación desde el cliente, facilitando así su uso en múltiples dispositivos y plataformas, sin las limitaciones de una aplicación de escritorio.
    
    \item Implementar mecanismos de seguridad avanzados. La aplicación incorporará diversas capas de seguridad, incluyendo autenticación de usuarios y cifrado robusto para proteger las credenciales. Se prestará especial atención a la prevención de vulnerabilidades, mediante el uso de sesiones seguras, códigos aleatorios y la gestión adecuada de estados.
    
    \item Optimizar la eficiencia y rendimiento de la aplicación. Se aplicarán técnicas y buenas prácticas para minimizar los tiempos de carga y garantizar que la aplicación funcione de manera rápida y eficiente, incluso cuando se maneje un gran volumen de datos.

    \item Generar un código limpio, modular y bien documentado. El código será diseñado para ser fácilmente reusable y mantenible, con una arquitectura desacoplada que permita modificar los inputs desde archivos externos. Además, se integrarán logs exhaustivos y documentación clara para facilitar futuras mejoras y el trabajo de otros desarrolladores.

    \item Aplicar diagramas de secuencia UML. Para garantizar un desarrollo estructurado y coherente, se emplearán diagramas de secuencia UML, permitiendo una planificación detallada de los flujos de trabajo y las interacciones entre los componentes del sistema.

\end{enumerate}
